\section*{Chapter \ref{ch:g11:modules-and-interfaces} --- Modules, Interfaces and Libraries}

\begin{solution}{exo:g11:modules-and-interfaces:1}
\code{geometry.distance((0, 0), (1, 1))} works and returns $\sqrt{2}$. \code{distance(...)} raises \code{NameError}:
\code{import geometry} makes only the name \code{geometry} available. \code{geometry.math.sqrt(4)} works and returns
\code{2.0}, but \code{math} is a \omterm{def:g11:modules-and-interfaces:module}{module} that \code{geometry} imported for itself, an \omterm{def:g11:modules-and-interfaces:interface}{implementation} detail: a user
imports \code{math} directly.
\end{solution}

\begin{solution}{exo:g11:modules-and-interfaces:2}
Running \code{main.py} prints \code{report}, when the import runs \code{report.py}, and then \code{__main__}.
Running \code{report.py} as a program prints \code{__main__}.
\end{solution}

\begin{solution}{exo:g11:modules-and-interfaces:3}
\code{MAX_SIZE}, \code{load} and \code{save}, provided they are documented. \code{_cache} and \code{_parse} are internal
by their underscore; \code{json} is a \omterm{def:g11:modules-and-interfaces:module}{module} imported for the \omterm{def:g11:modules-and-interfaces:module}{module}'s own use.
\end{solution}

\begin{solution}{exo:g11:modules-and-interfaces:4}
Sorted, the values are $1, 3, 7, 9$; an even count gives the mean of the two middle ones: \code{5.0}.
\texttt{[(\textquotesingle{}a\textquotesingle{}, 3), (\textquotesingle{}n\textquotesingle{}, 2)]}: three \code{a}, two \code{n}, one \code{b}.
\end{solution}

\begin{solution}{exo:g11:modules-and-interfaces:5}
In \code{temperature.py}:

\smallskip\noindent\hspace*{2em}\code{def to_fahrenheit(c):}

\noindent\hspace*{4em}\code{"""Return the temperature c, in degrees Celsius, in degrees Fahrenheit."""}

\noindent\hspace*{4em}\code{return c * 9 / 5 + 32}

\smallskip\noindent\hspace*{2em}\code{def to_celsius(f):}

\noindent\hspace*{4em}\code{"""Return the temperature f, in degrees Fahrenheit, in degrees Celsius."""}

\noindent\hspace*{4em}\code{return (f - 32) * 5 / 9}

\smallskip\noindent The tests check $100$ against $212$ and that $-40$ is the same on both scales.
\end{solution}

\begin{solution}{exo:g11:modules-and-interfaces:6}
\code{python stops.py} runs \code{stops.py} as \code{__main__}. Its first line starts loading \code{lines}, whose first
line imports \code{stops}: a second \omterm{def:g11:modules-and-interfaces:module}{module}, made from the same file. That one imports \code{lines}, which is already
being loaded and is returned as it is, without \code{LINES} yet; it then defines \code{FIRST_STOP} and reads
\code{lines.LINES}: the error. Two \omterm{def:g11:modules-and-interfaces:module}{modules} that each need the other at import time form a cycle. Correction: put
\code{FIRST_STOP} in a third \omterm{def:g11:modules-and-interfaces:module}{module}, \code{places.py}, which both import, and which imports neither.
\end{solution}

\begin{solution}{exo:g11:modules-and-interfaces:7}
The user's: the \omterm{def:g11:modules-and-interfaces:interface}{interface} never promised a list, and \code{sort} and \code{[0]} rely on the \omterm{def:g11:modules-and-interfaces:interface}{implementation}. The set
has no \code{sort} and no index. The right way is to ask the \omterm{def:g11:modules-and-interfaces:module}{module}'s author for a function in the \omterm{def:g11:modules-and-interfaces:interface}{interface}, such as
\code{earliest(s)}, implemented by each version (\code{min(s)} for the list and the set, \code{s[0]} for the sorted
list).
\end{solution}

\begin{solution}{exo:g11:modules-and-interfaces:8}
\smallskip\noindent\hspace*{2em}\code{def perimeter(points):}

\noindent\hspace*{4em}\code{pairs = zip(points, points[1:] + points[:1])}

\noindent\hspace*{4em}\code{return sum(distance(p, q) for p, q in pairs)}

\smallskip\noindent \code{points[1:] + points[:1]} is the list shifted by one place, the first point moved to the end:
\code{zip} pairs each point with the next, and the last with the first. The tests compare it with the original on
$100$ random polygons.
\end{solution}

\begin{solution}{exo:g11:modules-and-interfaces:9}
In \code{dateset_list}, \code{contains} compares with up to $n$ dates, and \code{add} too, since it first checks that
the date is absent. In \code{dateset_set}, both compute where the date is stored, and look at about one date,
whatever $n$ (chapter \ref{ch:g11:cost-of-an-algorithm}).
\end{solution}

\begin{solution}{exo:g11:modules-and-interfaces:10}
\code{1024}, then \code{1024.0}. \code{from math import *} copies all of \code{math}'s public names, among them
\code{math.pow}, which replaces the built-in \code{pow}; \code{math.pow} converts its arguments to \code{float}. The
import changed the meaning of a name the program already used.
% ledger: g11:py-math-pow
\end{solution}

\begin{solution}{exo:g11:modules-and-interfaces:11}
\smallskip\noindent\hspace*{2em}\code{def word_count(text, n):}

\noindent\hspace*{4em}\code{words = [w.strip(".,;:!?") for w in text.lower().split()]}

\noindent\hspace*{4em}\code{return Counter(w for w in words if w).most_common(n)}

\smallskip\noindent On \code{"The cat saw the dog. The dog saw the cat run."} with $n = 3$:
\texttt{[(\textquotesingle{}the\textquotesingle{}, 4), (\textquotesingle{}cat\textquotesingle{}, 2), (\textquotesingle{}saw\textquotesingle{}, 2)]}; \code{cat} comes before \code{saw} and \code{dog}, met first.
\end{solution}

\begin{solution}{exo:g11:modules-and-interfaces:12}
One possible \omterm{def:g11:modules-and-interfaces:interface}{interface}:
\begin{itemize}
  \item \code{new_queue()}: returns a new empty queue;
  \item \code{push(q, x)}: adds \code{x} at the end of \code{q};
  \item \code{pop(q)}: removes and returns the element at the front, the earliest pushed of those still there; on an
    empty queue it raises \code{IndexError} and leaves the queue unchanged;
  \item \code{is_empty(q)}: returns \code{True} when \code{q} holds no element;
  \item \code{size(q)}: returns the number of elements of \code{q}.
\end{itemize}
Nothing says how the elements are kept: a list, or two lists, could implement it, and a later chapter does.
\end{solution}

\begin{solution}{exo:g11:modules-and-interfaces:13}
\smallskip\noindent\hspace*{2em}\code{import os}

\noindent\hspace*{2em}\code{if os.environ.get("DATESET", "set") == "list":}

\noindent\hspace*{4em}\code{from dateset_list import add, contains, empty, from_dates, size}

\noindent\hspace*{2em}\code{else:}

\noindent\hspace*{4em}\code{from dateset_set import add, contains, empty, from_dates, size}

\smallskip\noindent The choice is made once, at the first import. A test file can run the same tests twice, in two
separate runs, \code{DATESET=list python tests.py} and \code{DATESET=set python tests.py}; or, within one run, set
the variable, remove \code{dates_chosen} from \code{sys.modules} so that it is imported again, and import it again, as
the chapter's tests do.
\end{solution}

\begin{solution}{exo:g11:modules-and-interfaces:14}
\omterm{def:g11:modules-and-interfaces:interface}{Implementation} changed, \omterm{def:g11:modules-and-interfaces:interface}{interface} kept: no line of the twenty programs changes, or at most the import line of each,
$20$ lines, if the \omterm{def:g11:modules-and-interfaces:module}{module}'s name changes. \omterm{def:g11:modules-and-interfaces:interface}{Interface} changed: each of the $20 \times 10 = 200$ calls is rewritten, in
programs written by others, and any missed one fails only when it runs. A change behind the \omterm{def:g11:modules-and-interfaces:interface}{interface} costs the
\omterm{def:g11:modules-and-interfaces:module}{module}'s author alone; a change of the \omterm{def:g11:modules-and-interfaces:interface}{interface} costs every user.
\end{solution}

\begin{solution}{exo:g11:modules-and-interfaces:15}
\smallskip\noindent\hspace*{2em}\code{LIMIT = 1_000_000}

\noindent\hspace*{2em}\code{_sieve = None}

\smallskip\noindent\hspace*{2em}\code{def is_prime(n):}

\noindent\hspace*{4em}\code{if _sieve is None:}

\noindent\hspace*{6em}\code{_build()}

\noindent\hspace*{4em}\code{return _sieve[n]}

\smallskip\noindent \code{_build} holds the sieve of the old top level, assigning the global \code{_sieve}; the message
is gone. Importing now only defines names; the first call to \code{is_prime} pays for the sieve, once. The tests check
that importing prints nothing and that both versions agree below $2000$.
\end{solution}

\begin{solution}{pb:g11:modules-and-interfaces:1}
\textbf{1.} $2028$: divisible by $4$, not a century year: $366$ days. $2100$: a century year not divisible by $400$:
$365$. $2000$: divisible by $400$: $366$.

\textbf{2.} $366$.

\textbf{3.} An impossible date is refused where it is made, with a message that names the problem, instead of
entering a set silently and causing a wrong answer far away.

\textbf{4.} \code{weekday()} is $4$: a Friday. $100$ days later is 28 February 2027.

\textbf{5.} \code{empty()} returns a new set with no date; \code{from_dates(dates)} returns a new set holding the dates
given, each once; \code{add(s, d)} adds \code{d} to \code{s}, and does nothing if it is there; \code{contains(s, d)}
tells whether \code{d} is in \code{s}; \code{size(s)} returns the number of dates in \code{s}.

\textbf{6.} \code{test_add_twice}.

\textbf{7.} So that they check only the promises, which every \omterm{def:g11:modules-and-interfaces:interface}{implementation} must keep; a test using \code{append}
would fail on a correct \omterm{def:g11:modules-and-interfaces:interface}{implementation} with a set.

\textbf{8.} \code{implementation: dateset_sorted}, then the three lines ending in \code{ok}, the line of dashes and
\code{OK}.

\textbf{9.}

\smallskip\noindent\hspace*{2em}\code{def test_from_nothing(self):}

\noindent\hspace*{4em}\code{self.assertEqual(ds.size(ds.from_dates([])), 0)}

\smallskip
\textbf{10.} \code{contains} relies on it: \code{bisect_left} halves the list as binary search does, which is right only
on a sorted list. \code{from_dates} sorts, and \code{add} inserts with \code{insort} at the right place.

\textbf{11.} The list: up to $100\,000$ dates. The sorted list: $h(10^5) + 1 = 17$.

\textbf{12.} \code{list(dates)} would keep repeated dates, and \code{size} would count them twice;
\code{dict.fromkeys} keeps each date once, in the order of first appearance.

\textbf{13.} After the search, \code{insort} must make room: every later date moves one place, up to $n$ moves.

\textbf{14.} From $5.72$ to $501.27$~$\mu$s: about $88$ times, for $100$ times more dates. The count predicts $100$;
at small $n$, the fixed cost of each call weighs more.

\textbf{15.} They hardly change: the sorted list stays near $0.2$~$\mu$s ($17$ dates looked at instead of $10$), the set
near $0.06$~$\mu$s.

\textbf{16.} About ten times more: $5$~ms per lookup.

\textbf{17.} The list, $1000 \times 501$~$\mu$s, about half a second; the sorted list, about $0.23$~ms; the set, about
$0.06$~ms.

\textbf{18.} Each \code{add} on the list scans it: $n$ calls look at about $n^2/2$ dates, $5 \times 10^9$ for $n =
10^5$, far too long for a measurement; \code{from_dates} is linear.

\textbf{19.} The import lines and the list \code{[dateset_list, dateset_set]}; none of the lines using the dates.

\textbf{20.} At $n = 10^5$ dates, \textbf{the set answers \code{contains} about $7900$ times faster than the list}
($0.06$ against $501$~$\mu$s per lookup, on the author's laptop), and \textbf{the same tests pass, unchanged, on every
\omterm{def:g11:modules-and-interfaces:interface}{implementation}}.
\end{solution}
